\begin{otherlanguage*}{english}
\noindent\textbf{Chapter XVII, manuscript pages 1--70: attachments, diagrams, and derived left Kan extensions}

\medskip
\noindent\emph{Commentary and audit produced by GPT-6 Astra (OpenAI, ChatGPT Pro mode). This credit concerns the present mathematical commentary and its accompanying registers; it does not certify any subsequent assembly or editing.}

\medskip
Grothendieck is looking for a way to construct the operations of a derivator from a category endowed with weak equivalences and cofibrations \citep{Grothendieck1990XVII}. The principal operation in these pages is a left adjoint to restriction of diagrams. Its construction has two logically separate parts. First one must find diagrams on which an ordinary colimit or ordinary left Kan extension respects weak equivalences. Then one must show that the functor obtained after localization has the required adjunction. Much of the force of the chapter lies in its refusal to identify these two tasks. The counterexample in Remark~3.6 makes the distinction precise, and the mapping-cone argument in Theorem~3.7 supplies the missing adjunction for the indexing posets under consideration.

The first twenty-four manuscript pages are a preliminary investigation which the author asks the reader to skip once section~2 has been written. They nevertheless disclose the mathematical problem that determines the stronger definition. Throughout, an open subset of an ordered set is a lower ideal, and $\mathcal M(I^\circ)$ denotes covariant $I$-diagrams. For a map $u:X\to Y$, the elementary test is that
\[
 X_i\amalg_{\mathop{\rm colim}_{j<i}X_j}
       \mathop{\rm colim}_{j<i}Y_j\longrightarrow Y_i
\]
be a cofibration. This relative latching map measures the attachment at $i$ after all its predecessors have been incorporated. If the indexing order is well founded, one can perform these attachments in a compatible ordinal order. Pushout stability and transfinite-composition stability then put both every component of $u$ and its colimit map in $C$. Applying the same argument to $TC=C\cap W$, and using two-out-of-three at the next attachment, identifies trivial diagram cofibrations with the diagram cofibrations that are objectwise weak equivalences. The commentary supplies a continuous filtration so that this argument does not tacitly assume continuity of the original diagrams.

The relative statement is more useful than the absolute one. For lower ideals $K\subset J$, write $X_K=\mathop{\rm colim}_{K}X$ and similarly for $Y$, and put
\[
 r_{K,J}(u):Y_K\amalg_{X_K}X_J\longrightarrow Y_J.
\]
The essential excision identity says that when a new lower ideal $B$ is adjoined to $A$, the relative map over $A\subset A\cup B$ is a pushout of the relative map over $A\cap B\subset B$. A succession of such comparisons is controlled by a continuous filtration with terms $X_J\amalg_{X_{A_\alpha}}Y_{A_\alpha}$. These two elementary colimit calculations account for the large diagrams of the preliminary section and for Lemmas~2.6--2.7. They also explain the distinction between the relative assertion, which does not require $X$ to be cofibrant, and the assertion $X_K\to X_J\in C$, which does. Footnote~5 overstates the dispensability of this latter hypothesis; footnote~7 already gives the necessary qualification.

Section~2 makes all these relative comparisons part of the definition. A map belongs to the new $C_I$ when $r_{K,I/i}(u)\in C$ for every $i$ and every lower ideal $K\subset I/i$. Theorem~2.1(1) extends that condition to every inclusion of lower ideals, without assuming that $I$ is artinian. Part~(2) recovers the single relative-latching test when $I$ is artinian. Part~(3) identifies $(TC)_I$ with $C_I\cap W_{I^\circ}$ under the same well-foundedness and closure hypotheses. This change of definition is indispensable: on a poset with an endless descending chain, all relative latching maps can be isomorphisms while the components of a map of diagrams fail even to be injective. The development attached to manuscript page~24 proves that this obstruction detects every failure of well-foundedness, using diagrams of sets extended from the upward closure of a descending chain.

The question left on manuscript pages~38--39 concerns part~(3), not part~(2): can the artinian hypothesis be removed from the equality involving $TC$ when the stronger definition of $C_I$ is retained? Grothendieck verifies one descending-chain example. The present commentary extends that argument to linearly ordered sets and, more generally, to posets whose principal lower ideals are linearly ordered. A well-ordered cofinal subset of a lower chain reduces its colimit to the established ordinal case. We did not find a resolution of the unrestricted question under exactly the assumptions of Theorem~2.1 in the literature inspected for this package. This is a statement about the scope of our verification, not a claim that the problem is known to remain open. The later theorems on diagram cofibration structures use their own definitions and additional hypotheses, and therefore do not by themselves settle this precise question.

Several corrections concern the formulation of an argument whose conclusion remains valid. Lemma~2.7 must use a chain continuous at limit ordinals; an arbitrary increasing chain can acquire new objects at a limit stage, so its successor comparisons cannot account for the entire relative map. The printed lemma also confuses its initial relative domain with the initial object of an absolute filtration. Both points are repaired explicitly. The description of the intermediate cube in Lemma~2.6 as totally cocartesian is too strong if it means that all six faces are pushouts; pushout associativity proves the needed face without that assertion. These repairs are recorded separately from transcription corrections. The edition's printed slips remain in the source text, while the commentary states and proves the mathematically intended claims.

Section~3 turns the attachment calculus into homotopical algebra. The axioms $\mathrm{Cof}_0$--$\mathrm{Cof}_3$ require two-out-of-three for $W$, composition and pushout stability for $C$, pushout stability for $TC$, and a factorization of every map as a cofibration followed by a weak equivalence. They do not require every object to be cofibrant. The full subcategory of cofibrant objects has the same localization, so Brown's methods can be applied there. The factorization of a diagram is then constructed one index at a time: a factorization already chosen on the predecessors of $i$ produces the map
\[
 X_i\amalg_{L_iX}L_iZ\longrightarrow Y_i,
\]
to which the ambient factorization axiom is applied. This is the direct-diagram form of the Reedy construction. Its later cellular formulation is made explicit in \citet[Proposition~6.3 and Lemma~7.1]{RiehlVerity2014Reedy}. The source works also with ordinal ranks, which exceed the natural-number degree convention used for many familiar Reedy categories.

The two finiteness branches in Theorem~3.1 deserve to be read literally. In the first, $I$ is artinian and its lower-Alexandrov topology is noetherian. Under the ordinary choice convention used in the argument, their conjunction forces $I$ to be finite: otherwise successive choices of a minimal element of the remaining poset give a strictly increasing sequence of finite lower ideals. The second branch allows arbitrary small artinian posets, requires the necessary small colimits, and imposes transfinite stability on both $C$ and $TC$. Thus the theorem has a finite construction and a construction with explicit infinitary hypotheses. The commentary keeps these alternatives separate when it discusses the available Kan extensions.

Proposition~3.2 proves that left Kan extension preserves diagram cofibrations, and that restriction preserves them under a suitable local cofilteredness condition. Brown's cylinder factorization then shows that left Kan extension carries weak equivalences between cofibrant diagrams to weak equivalences. Together with the localization equivalence for cofibrant objects, this yields the candidate $u_!^W$. The warning in Remark~3.6 is categorical and complete: choosing a subcategory which computes a localization need not make its inclusion a derived replacement suitable for an adjunction. With the identity adjunction on the walking arrow, different choices of weak equivalences give the stated failure. The related assertion that full faithfulness survives localization also fails without control of the localized counit. These are failures of general inferences, not counterexamples to the properly constructed derived adjunction.

Theorem~3.7 restores the adjunction by factoring $u:I\to I'$ through its categorical mapping cone. One factor is a closed inclusion, whose left Kan extension is extension by the initial object. The other is a retraction whose relevant comma categories have initial objects. Each factor consequently falls under a case where the adjunction can be checked, and their left adjoints compose. The discussion of arbitrary target categories that follows needs a further size qualification in the finite branch. For $u:e\to B\mathbb Z$ and $\mathcal M=\mathrm{FinSet}$, a left adjoint to the forgetful functor from finite $\mathbb Z$-sets would have to provide a free finite $\mathbb Z$-set on one generator; finite cycles of unbounded order show that no such object exists. A finite source does not make its comma categories finite. This obstruction disappears when the required small coproducts and colimits are supplied, but that hypothesis must be stated.

The chapter belongs to the program described in \emph{Pursuing Stacks}, section~69, and in the letter to Thomason of 2~April~1991: diagram localizations retain operations that a single homotopy category does not specify, while cofibrations and fibrations are auxiliary means of constructing those operations \citep{Grothendieck1983,Grothendieck1991Letter}. Later work establishes the desired diagram formalism with precise hypotheses. \citet[Corollary~2.24]{Cisinski2010} obtains a one-sided derivator on finite direct categories: in his convention, a weak right derivator satisfying the additional axiom Der~5. For ABC cofibration categories, \citet[Theorems~9.6.3, 9.6.5, and 10.3.2]{RadulescuBanu2009} gives arbitrary small homotopy left Kan extensions, their adjunctions, pointwise formulas, and the resulting derivator. \citet[Theorem~3.33]{Lenz2018} identifies the homotopy prederivator of an ABC cofibration category in which every object is cofibrant with that of its quasicategory of frames. These comparisons place the construction in modern homotopy theory while preserving the distinctions in assumptions that matter in this chapter.

The last paragraph introduces strictness in order to apply the same construction to categories of diagrams with an arbitrary additional parameter. Functorial factorization gives objectwise diagram factorizations; the universal arrow diagram gives the converse when the relevant arrow category is small. Here the source returns explicitly to its historical presheaf convention, $\mathbf D_W(J)=\operatorname{Fun}(J^{\mathrm{op}},\mathcal M)[W^{-1}]$. The preceding covariant construction concerns $H(I)=\mathbf D_W(I^{\mathrm{op}})$, so its immediate infinite range must be translated by opposition; the finite range is unchanged. In this presheaf convention the final printed comma category $i'\backslash I=(i'\downarrow u)$ is correct: the colimit computing a left Kan value runs over its opposite. The covariant formula instead uses $(u\downarrow i')=I/i'$. A walking-arrow calculation in both conventions explains the difference and prevents a change of variance from being misdiagnosed as an edition slip. The later ABC theorem supplies the arbitrary-shape extension under the full infinitary hypotheses. The manuscript's own further construction and forward references lie beyond this excerpt.
\end{otherlanguage*}
